% =============================================================================
% first_document.tex
% Workshop: Scientific Writing with LaTeX, Overleaf & Prism
% Purpose:  A complete beginner LaTeX document demonstrating core features.
%           This file compiles without errors as-is.
% =============================================================================

% --- Document class -----------------------------------------------------------
% "article" is the most common class for short documents, papers, and homework.
% Options: 12pt sets the base font size; a4paper sets the page dimensions.
\documentclass[12pt, a4paper]{article}

% --- Packages -----------------------------------------------------------------
% Each \usepackage line loads additional functionality.

\usepackage[utf8]{inputenc}      % Allow UTF-8 characters in the source
\usepackage[T1]{fontenc}         % Better font encoding for accented characters
\usepackage{amsmath, amssymb}    % Extended math support
\usepackage{graphicx}            % Include images with \includegraphics
\usepackage{booktabs}            % Professional-quality tables (\toprule, etc.)
\usepackage{hyperref}            % Clickable cross-references and URLs
\usepackage{geometry}            % Easy page-margin control
\geometry{margin=2.5cm}

% --- Title block --------------------------------------------------------------
\title{My First \LaTeX{} Document}
\author{Your Name Here}
\date{\today}   % \today inserts the current date automatically

% ==============================================================================
\begin{document}

% Generate the title from the \title, \author, and \date above
\maketitle

% --- Table of contents (optional for short docs) -----------------------------
\tableofcontents
\newpage

% ==============================================================================
% SECTION 1 — Introduction
% ==============================================================================
\section{Introduction}
\label{sec:introduction}

This document demonstrates the most common features of \LaTeX{}.
It was created as part of the \textbf{Scientific Writing with LaTeX, Overleaf
\& Prism} workshop.

You can reference other parts of the document automatically.
For example, the table appears in Section~\ref{sec:tables} and the equation
is in Section~\ref{sec:equations}.

% ==============================================================================
% SECTION 2 — Text Formatting
% ==============================================================================
\section{Text Formatting}
\label{sec:formatting}

% --- Bold, italic, and other emphasis ----------------------------------------
\subsection{Emphasis and Font Styles}

Here are the most common ways to style text:

\begin{itemize}
    \item \textbf{Bold text} using \verb|\textbf{...}|
    \item \textit{Italic text} using \verb|\textit{...}|
    \item \texttt{Monospace / code} using \verb|\texttt{...}|
    \item \underline{Underlined text} using \verb|\underline{...}|
    \item \emph{Emphasised text} using \verb|\emph{...}| — adapts to context
\end{itemize}

% --- Lists --------------------------------------------------------------------
\subsection{Lists}

\textbf{Numbered list} (enumerate):
\begin{enumerate}
    \item Write your content in a \texttt{.tex} file.
    \item Compile with \texttt{pdflatex} (or use Overleaf).
    \item Review the resulting PDF.
\end{enumerate}

\textbf{Bullet list} (itemize):
\begin{itemize}
    \item First item
    \item Second item
        \begin{itemize}
            \item A nested sub-item
            \item Another sub-item
        \end{itemize}
    \item Third item
\end{itemize}

% ==============================================================================
% SECTION 3 — Equations
% ==============================================================================
\section{A Simple Equation}
\label{sec:equations}

The quadratic formula gives the solutions of $ax^2 + bx + c = 0$ as:

\begin{equation}
    x = \frac{-b \pm \sqrt{b^2 - 4ac}}{2a}
    \label{eq:quadratic}
\end{equation}

We can refer to it as Equation~\eqref{eq:quadratic}.
Inline math uses single dollar signs: $E = mc^{2}$.

% ==============================================================================
% SECTION 4 — Figures
% ==============================================================================
\section{Including a Figure}
\label{sec:figures}

% The figure below is ready to use — just uncomment \includegraphics and
% replace "example-image" with your filename (without extension).

\begin{figure}[htbp]          % h=here, t=top, b=bottom, p=own page
    \centering
    % Uncomment the next line and provide your image file:
    % \includegraphics[width=0.6\textwidth]{example-image}
    \rule{0.6\textwidth}{4cm}  % Placeholder rectangle — remove when using a real image
    \caption{A placeholder figure. Replace with your own image.}
    \label{fig:placeholder}
\end{figure}

As shown in Figure~\ref{fig:placeholder}, you can reference figures by label.

% ==============================================================================
% SECTION 5 — Tables
% ==============================================================================
\section{Tables}
\label{sec:tables}

Table~\ref{tab:example} shows a simple data table using the \texttt{booktabs}
package for professional horizontal rules.

\begin{table}[htbp]
    \centering
    \caption{Experimental results summary.}
    \label{tab:example}
    \begin{tabular}{lcc}
        \toprule
        \textbf{Method} & \textbf{Accuracy (\%)} & \textbf{Time (s)} \\
        \midrule
        Baseline         & 72.3  & 1.2 \\
        Proposed (v1)    & 85.1  & 2.4 \\
        Proposed (v2)    & 89.7  & 3.1 \\
        \bottomrule
    \end{tabular}
\end{table}

% ==============================================================================
% SECTION 6 — Cross-References Summary
% ==============================================================================
\section{Cross-References}
\label{sec:crossref}

\LaTeX{} resolves references automatically when you compile twice:
\begin{itemize}
    \item Section~\ref{sec:introduction} — Introduction
    \item Equation~\eqref{eq:quadratic} — Quadratic formula
    \item Figure~\ref{fig:placeholder} — Placeholder figure
    \item Table~\ref{tab:example} — Results table
\end{itemize}

% ==============================================================================
% SECTION 7 — Conclusion
% ==============================================================================
\section{Conclusion}

You now have a working \LaTeX{} document!  Try modifying each section to
practise what you have learned.

\end{document}
